\section{Introduction}\label{sec:intro}

The Rule of 40 holds that a software company's revenue growth rate plus its profit margin should reach 40\%. It is a widely used benchmark in software investing \citep{feld2015rule,depeyrot2018hacking}, and it has spawned variants: a ``Rule of 60'' for AI-enabled software \citep{vista2026rethinking}, growth-weighted versions such as the ``Rule of X'' \citep{bessemer2024rulex}, and the suggestion that firms reinvesting in AI may for a time have to settle for a ``Rule of 30'' \citep{lipman2026ai}. Reported scores at the top have reached far above 40: Palantir reported a score of 155 for the second quarter of 2026 on its adjusted margin \citep{palantir2026q2}. The practitioner discussion does not say which combinations of growth and margin a firm can attain, what the frontier of attainable combinations looks like, whether 40 is a point on that frontier or a line that crosses it, or how high the score could go if artificial intelligence automates most of the work a software company does.

This paper asks four questions. What is the efficient frontier of growth and profitability, and how does the Rule of 40 relate to it? What is the theoretical limit of the score: is it a Rule of 100, a Rule of 1000, or something else? How might that limit evolve over the next two, five and ten years as AI automates more of the firm? And what conditions must a firm meet to approach it?

We answer with a model, an empirical frontier and a calibration. In the model (Section~\ref{sec:model}), a firm grows faster than its market only by spending, and the returns to spending diminish with the amount spent and with the share of the market the firm already holds. The efficient frontier of attainable (growth, margin) pairs is then decreasing and concave, with an intercept equal to the margin the firm earns when it simply grows with its market. The Rule of 40 and the Rule of X are price lines, not frontiers: their slope is the number of margin points the rule trades for one point of growth, the relative price of growth. They are lines of equal score, not of equal value; firms on the same line can be worth very different multiples of revenue (from \IsoZero{} to \IsoThirty{} times revenue at a score of 40 in the illustration of Table~\ref{tab:isoscore} of the Supplementary Appendix), and this paper estimates no structural mapping from scores to multiples. The maximal score is reached where one of these lines is tangent to the frontier. In closed form it is the margin intercept, plus market growth net of the capital it requires, plus a growth surplus that vanishes as the firm's share approaches its cap. Within one market and at given retention, every point of that surplus is paid for out of a finite budget of share headroom, so extreme scores are necessarily transient. The model separates two weights on growth that practice treats as one: the weight investors use to screen firms, which reflects the duration of the revenue stream and is consistent with the weights of two to three in use, and the weight a value-maximizing firm acts on. Under a standard verification hypothesis (\autoref{prop:dynamic}), the latter is the marginal value of a dollar of revenue, which exceeds the screening weight when the firm's revenue multiple is well above it; how far above depends on the owner's discount rate. In value terms the Rule of X is a rule of economic profit plus growth at the net value of growth, and the cash Rule of 40 weights growth at one minus capital intensity, so that for capital intensity above one it penalizes growth whatever its return (\autoref{prop:general}); \autoref{rem:general-phi} lists the results that carry beyond software and the conditions they need. For a firm whose decision weight is highest at small share, fast growth gains most from automating the growth engine and retention, and maturity from automating overhead, while retention still adds to the durable score near saturation. At observed positions the ranking depends on how much selling spend buys growth: net of maintenance selling spend, cutting general and administrative cost outranks automating the growth engine at \lpNGACutFirstNet{} of \lpNPositions{} positions at the central maintenance level, and the growth engine leads at \StLeverMin{} to \StLeverMax{} of them over that level's bootstrap interval (Section~\ref{sec:levers-position}). Under free entry into a market for a homogeneous product, the long-run score is pinned down by market growth, the cost of capital, a normal return on customer capital and the rent on whatever input rivals cannot replicate. And the frontier estimated across firms is drawn through the points firms chose, so it does not identify the trade-off any one firm faces.

Two extensions bring the costs and the pace of automation into the frontier. First, full automation replaces labor with compute, fees, an accountability floor and the recurring costs of maintaining, securing and supervising AI. We give each floor a closed form. The accountability floor is a pass-fail control standard: the cheapest design that leaves no reasonable possibility of an undetected material misstatement. Control layers that share one AI model share its failures, so no design built only from such layers passes once their common-mode failure is more than remote, however many layers it stacks, and neither cheaper inference nor scale removes that requirement (\autoref{prop:amin}). Second, a firm reaches its full-cost frontier only through transitions that fail, roll back and draw on limited capacities to absorb change, decide and review deployments, and that new leaders, strategies and reorganizations can disrupt. In the frontier attainable by a given date, risk slows the approach to the static frontier, approximately like a cut in the program's pace, without lowering the frontier, while partial realization of benefits, when the realized share is a property of the process, lowers the level to which the firm converges (\autoref{prop:attainable} and \autoref{prop:expected}).

The empirical analysis (Sections~\ref{sec:data} and~\ref{sec:empirics}) uses SEC XBRL filings, taken as first reported, for \VFirmsEone{} prepackaged-software firms over fiscal years 2010 to 2025, including firms later acquired or delisted, after sample rules fixed before any estimate on the resulting sample. Because both growth and margin are desirable, we use partial-frontier statistics whose meaning does not depend on sample size. Observed scores are the baseline; a claim that confronts the model is stated only where it also holds on the organic frontier, which removes acquired growth, and under two imputation bounds for the years that removal drops. Six findings stand out. First, among the best performers a point of growth has come with well under a point less margin: over 2010--25 the slope of the 0.90-quantile frontier is $-\PsiSlopeOp$ on the GAAP basis and $-\PsiSlopeCash$ after stock-based compensation in the baseline, and \ScPsiOpPLo{} to \ScPsiOpPHi{} on the GAAP basis for 90\% of the defensible analysis choices of a specification curve, far from the one-for-one trade of the Rule-of-40 line in every view and every specification. These are points firms chose across firms, not any firm's own trade-off, and the negative sign does not survive a familywise correction. On the GAAP basis and for growth of up to 40\%, the slope was \PsiEarlyOp{} in 2010--19 and \PsiLateOp{} in 2020--25, a steepening negative in every view with an interval that excludes zero in \PsiDiffOpCountNeg{} views, but fragile: across the sample choices of the specification curve that interval excludes zero in only \ScDiffOpBootNinetySplitAShare\% of the samples tested, against \ScDiffOpBootNinetySplitBShare\% with a 2010--15 against 2018--25 split. After SBC the slope did not steepen in the samples that remove acquired growth. Within firms, margin and growth are nearly unrelated. Second, the data identify neither the curvature of the growth technology nor the weight on growth that firms act on: the ratio of growth spend to excess growth lies between \SpendLo{} and \SpendHi{} in the four views (\ScSpendMin{} to \ScSpendMax{} across sample choices), but the product it is meant to reveal ranges from \AlphaAllLo{} to \AlphaAllHi{} across methods. Third, the decline of the upper tail of the score with firm size is weak and not robust. Fourth, the lower cash-basis scores of capital-intensive firms are an accounting effect of deducting capital spending. Fifth, between 2016--19 and 2023--25 the 95th percentile of the score fell in every sample variant and every specification that measures it, by \TopFallOpEone{} points on the GAAP basis and \TopFallCashEone{} after stock-based compensation in the baseline, near the top of the ranges of \ScTopOpCoreFallLo{} to \ScTopOpCoreFallHi{} and \ScTopCashCoreFallLo{} to \ScTopCashCoreFallHi{} points that the sample choices give; the GAAP decline survives a familywise correction at 10\% under \TopFamTen{} of four procedures and at 5\% under \TopFamFive, and the cash decline under none. Fewer firms grew fast in 2023--25, but a reweighting by growth bins does not separate the change in the growth mix from changes within bins robustly, every quantile of the score fell, and we do not identify why the top fell. On the GAAP basis the 95th-percentile margin of firms growing at least 20\% changed by \TtwentyOp{} points, about as much as the score, with an interval that excludes zero under three of four interval types, and after SBC that of all firms with nonnegative growth changed by \TzeroCash{} points, with an interval that excludes zero in two of four views. The top partly recovered in 2025, by \RecoveryEone{} points under observed scores but only \RecoveryNoAcq{} without acquisition-affected firm-years. Sixth, firms near the frontier earn about \GapV{} points more operating margin than typical firms at the same growth rate; among firms growing 40 to 60\% a year, \GapVFastSM{} points of that band's \GapVFast-point gap come from lower selling and marketing expense.

Combining model and estimates (Section~\ref{sec:horizons}), we find that automation alone does not set the limit of the rule. Inverting the model firm by firm for the best scores of 2023--25 shows that they exceeded their durable level if little of their selling and marketing expense buys growth; how much of them was transient the data do not determine, and the anchor rests on \AncNEone{} firm-years of \AncFirmsEone{} firms. Under full automation, and holding prices fixed, the cash-basis score of a mature firm is bounded by its after-tax margin net of compute costs, fees, an accountability floor and the recurring costs of maintaining, securing and supervising AI, plus market growth net of capital needs. If compute absorbs \CcCen\% of revenue at full automation at 2025 prices, a scenario value, that bound is \CeilCenFive{} to \CeilCenThirty{} points for market growth of 5 to 30\%, so an asset-light firm exceeds 100 only if its market grows faster than about \GstarCenOne\% a year at the central cost inputs (\GstarFav\% with favorable parameters, and \StGstarCcSixMin{} to \StGstarCcSixMax\% with compute at \CcCen\% and the other costs moved jointly over their stated ranges). Every threshold is conditional on the scenario compute share: over its scenario range of \CcLo{} to \CcHi\% of revenue the central threshold runs from \GstarCcLo{} to \GstarCcHi\%. Along a scenario path of the compute share it is \GstarCenTwentyEight, \GstarCenThirtyOne{} and \GstarCenThirtySix\% in 2028, 2031 and 2036, and if the industry passes half of every unit-cost saving to customers, or 70\%, it rises to \GstarPhiHalfOne{} or \GstarPhiSevenOne\% at 2025 prices. At the market growth rates software markets have sustained, a Rule of 1000 is not a steady state: it requires \PullMultLo{} to \PullMultHi{} times the demand pull of today's best public software firms, and the share-headroom budget limits how long any such surplus can last. The pace of approach is limited by execution. Under central execution a top-tail firm (at the 95th percentile of 2023--25 scores) realizes about \PaceCenSix\% of the benefits of its automation program by 2036 (\StRealMin{} to \StRealMax\% over the stated ranges of execution), having embedded about \CenEmbSix\% of it, and under fast execution the upper bound on its embedded share reaches all of it in \PaceFastDone; while it is transforming, the costs of transitions in flight lower its cash score by about \CenDipEight{} to \CenDipSix{} points. On the central scenario path, which is not a forecast and holds output prices fixed, the maximal durable score of such a firm is about \CenREight{} in 2028, \CenROne{} in 2031 and \CenRSix{} in 2036 (about \CenRSixNet{} while it still pays for transitions in flight; 90\% anchor-sampling interval \CenRSixLo{} to \CenRSixHi, which holds the program and the target fixed, against \StNetSixPFive{} to \StNetSixPNinetyFive{} for 90\% of draws with every scenario input moved, the draws uniform over the stated ranges, log-uniform for compute per dollar of labor; Section~\ref{supp:stress}), when selling and marketing expense above a maintenance level of \MaintBase\% of revenue buys growth; if all of it does, the 2036 score is \CenRSixAllSM, and if none does, \CenRSixNoSM. By 2036 the largest swing in a one-at-a-time sensitivity comes from \TornFirstSix. Under central execution annual gains first exceed the costs of transitions in flight after about \PayCrossCen{} years, but the program repays those costs only after \PayBackCen{} years (\PayBackDiscCen{} at a \acqRwAnnual{} cost of capital; \StPayCenMin{} to \StPayCenMax{} years as the other inputs move jointly); at three and at five years its cumulative effect is still negative at all but a few points of that joint range, so at the central inputs the program fails a five-year cash-payback test, with exceptions only where compute starts high and falls steeply, although an owner who can sell at a price that reflects the remaining benefits gains from it. Under free entry a rent on a scarce input cannot lift the score above the full-cost ceiling (\autoref{prop:entry}(d)). The levers that move the limit are therefore a non-replicable input, exceptional demand pull, capital that earns its keep, pricing that captures falling unit costs, small accountability and full-cost floors and a fast-growing market; a sustained Rule of 100 requires them jointly to push the durable score above 100.

Measurement matters as much as economics at the top of the distribution. Adjusted margins that exclude stock-based compensation, capitalized computing capacity, gross revenue recognition and run-rate revenue can each push a reported score above the ceiling without any change in economics. Growth that is bought rather than built can be the largest such distortion: acquired revenue enters the growth term, while its price, paid in investing cash flow, never touches a cash score, and acquisition debt moves interest below the operating margin. Charged properly, acquired growth enters the score at the net weight $X-k^{\rm acq}$, where $k^{\rm acq}$ is the price paid per dollar of acquired revenue, and each acquisition in a market the firm does not already serve buys a fresh budget of share headroom\IfFileExists{sections/measurement_acquired.tex}{ (Section~\ref{sec:acquired-growth})}{}. Bending Spoons, a serial acquirer that listed in 2026, does not report a Rule-of-40 score. Computed from its reported revenue growth and GAAP operating margin, its score for the second quarter of 2026 was about \acqRQTwoRep, and on the company's own unaudited measure of organic growth the organic score runs from about \acqRQTwoOrgIntR{} to \acqRQTwoOrgNoReorgR, depending on how interest and reorganization expense are treated (on its annual figures for 2025, about \acqROrgIntTwentyFiveR{} to \acqROrgNoReorgTwentyFiveR). Its 2025 GAAP score, recomputed from its Form F-1, is \acqRGaapTwentyFive{} before interest and \acqRIntTwentyFive{} after. On a different basis, its cash margin before interest, and with its acquisitions charged as capital, so that financing is not charged twice, the 2025 score lies between \acqConvLo{} and \acqConvHi, depending on how the charge is spread over time and on a scenario life of the revenue bought. The company reports material weaknesses in its internal control over financial reporting, and we use its figures as filed. Acquired growth is why a claim about the frontier is stated only where it holds both in observed scores and on the organic frontier.

The paper builds on several literatures (Section~\ref{sec:lit}). The tangency between a growth--profit transformation curve and a valuation line goes back to \citet{marris1964economic}, \citet{solow1971some}, \citet{wood1975theory} and \citet{odagiri1981theory}; our growth technology is a reduced form of customer capital and market penetration \citep{arkolakis2010market,gourio2014customer} and resembles the advertising models of \citet{vidale1957operations}, \citet{nerlove1962optimal} and \citet{sethi1973optimal}, and recent evidence that firms grow by acquiring customers at roughly constant markups \citep{fitzgerald2024how,argente2025how} supports it. The peer-reviewed academic studies of the rule we are aware of evaluate it as an investment screen \citep{lee2024evaluating} or relate it to enterprise value \citep{nei2026enterprise}; neither estimates the frontier. Empirically, we adapt partial-frontier estimators \citep{cazals2002nonparametric,aragon2005nonparametric} to a frontier on which both outputs are desirable, and we map the effects of AI automation \citep{aghion2019artificial,acemoglu2025simple,babina2024artificial} onto the parameters that determine a firm's score. The results we believe are new are the score budget, the firm's decision weight on growth and its relation to screening weights, a worked instance, for growth--margin frontiers, of the classical non-identification of technologies from chosen points, the accountability floor as a control standard under common-mode failure, and the attainable and expected frontiers; the rest follows from the functional form or is a special case of a standard model. A companion paper \citep{murthy2026transformation} derives the model's parameters from the tasks a firm automates; this paper does not depend on it.

The author directed the research and made its decisions; AI agents drafted the text, code and proofs (Appendix~\ref{app:protocol}). The appendix records which agents contributed and on which models, the review rounds, the decisions the author made, the choices made after results were seen and what the protocol does not establish; the Supplementary Appendix gives the detailed records.

Section~\ref{sec:lit} reviews related work. Section~\ref{sec:model} presents the model. Section~\ref{sec:data} describes the data and Section~\ref{sec:empirics} the empirical frontier.\IfFileExists{sections/measurement_acquired.tex}{ Section~\ref{sec:acquired-growth} shows how to measure the score when growth is bought.}{} Section~\ref{sec:horizons} calibrates the ceiling and a central path by horizon, Section~\ref{sec:discussion} discusses the weights on growth, measurement and the accountability floor, and Section~\ref{sec:conclusion} concludes. The Supplementary Appendix, a separate document, contains the proofs, the sources of every parameter, further robustness tables and the detailed protocol records.

\paragraph{How not to use these numbers.} The results are easy to quote out of context, so we state the main misreadings here.
\begin{enumerate}[leftmargin=*]
\item The ceiling of Section~\ref{sec:horizons-ceiling} is an upper bound at a stated full-automation cost structure and fixed prices. It is not a forecast of scores, not a target and not a cap on valuation, and a firm whose costs are leaner than that structure can exceed it.
\item The threshold of \GstarCenOne\% holds for a firm with no capital needs, compute at \CcCen\% of revenue at 2025 prices and the other costs at their central values (\StGstarCcSixMin{} to \StGstarCcSixMax\% as those costs move jointly over their stated ranges), which the path of the compute share reaches only in about \CcYearSix. It is lower for a firm whose billing runs ahead of its revenue, so that its cash capital intensity is negative (about \GstarKdefMedCen\% at the median deferred-revenue intensity of growing firms), and higher if savings pass to customers.
\item The central path is a scenario, not a forecast; quote it with its interval and bounds. In 2036 it is \CenRSix{} (\CenRSixLo{} to \CenRSixHi{} from anchor sampling alone, \StNetSixPFive{} to \StNetSixPNinetyFive{} for 90\% of draws with every scenario input moved over its stated range, uniformly or, for compute per dollar of labor, log-uniformly, an assumption, not a measured distribution), \CenRSixAllSM{} if all selling and marketing expense buys growth and \CenRSixNoSM{} if none does. The transient-surplus rows of Table~\ref{tab:tornado} are an order of magnitude for one firm type, not a path.
\item The slopes of Section~\ref{sec:empirics} describe the points firms chose, not any firm's own trade-off, and the fall of the top of the score distribution since 2016--19 is a description, not a case for lowering the bar or an account of why it fell. Their baseline values are one specification among many defensible ones: quote the slope and the fall with the ranges of the specification curve (Section~\ref{supp:speccurve}), and the steepening of the slope since 2020 as fragile.
\item The data do not identify the weight firms act on. A firm that cites a decision weight above the screening weights to justify a lower score must show, from the payback and retention of its customer cohorts at a stated discount rate, that its marginal growth spend earns more than it costs.
\item The lever ranking of Section~\ref{sec:levers-position} compares two levers at each position's own weight on a stated basis of growth spend; it is not a ranking of every lever or a plan for a particular firm.
\item The remote threshold of the control model is a scenario value, not a tolerance for misstatement, and the cheapest design in its example is not a control-design recommendation: admissibility is decided by management's assessment and the auditor's evidence.
\item Scores for named companies are our computations from their filings under stated assumptions, not the companies' measures, and scores with acquisitions charged in full are lower bounds in years whose outlay exceeds the annuity on earlier outlays still within their life, as in a serial acquirer's buying years (\autoref{rem:acquired-bounds}).
\item Quote every score with its margin basis, its growth measure (reported or organic) and whether interest is deducted.
\item Full automation is an analytical bound, not a staffing plan, and nothing here implies that accountable people or control owners can be removed.
\end{enumerate}
